\documentclass[aspectratio=169,11pt]{biscuitsbeamer}

\showexamplenumberfalse   % 取消例题编号
\showexercisenumberfalse  % 取消练习题编号
\title{函数的奇偶性}
\subtitle{高中数学·必修第一册}
\author{Biscuits}
\date{2025年9月20日}

\begin{document}

% ========== 封面 ==========
{
	\setbeamertemplate{footline}{}
	\begin{frame}[plain,noframenumbering]
		\TitlePageArtwork
		\vspace*{1.6cm}
		\hspace*{1.25cm}
		\begin{minipage}[t]{0.50\paperwidth}
			\vspace{0.5cm}
			{\color{PrimaryGreen}\bfseries\fontsize{24}{30}\selectfont 函数的奇偶性\par}
			\vspace{0.2cm}
			{\color{WarmGray}\Large 高中数学·必修第一册\par}
			\vspace{0.2cm}
			{\color{WarmGray}\normalsize 第一课时\par}
			\vspace{0.55cm}
			{\Large Biscuits\par}
			\vspace{0.2cm}
			{\large 2025年9月20日\par}
		\end{minipage}
	\end{frame}
}
% ========== 学习目标 ==========
\begin{frame}{本节学习目标}
	\begin{quan}
		\item 理解函数奇偶性的定义，能判断简单函数的奇偶性。
		\item 掌握奇函数、偶函数的图象特征。
		\item 掌握奇函数在 $x=0$ 处的性质。
		\item 能利用奇偶性解决求解析式、求值等问题。
	\end{quan}
	\vspace{3mm}
	\begin{tip}
		重点：奇偶性的定义与判断方法。\\
		难点：利用奇偶性求解析式、解决综合问题。
	\end{tip}
\end{frame}

% ========== 目录 ==========
\begin{frame}[allowdisplaybreaks]{目录}
	\tableofcontents
\end{frame}







\input{全国1卷.tex}








\end{document}